\begin{trapbox}
\begin{minipage}{\linewidth}
\textbf{\textcolor{CTrap}{易错一｜把“极线”一律理解成“切点弦”}}

\begin{center}
\begin{tikzpicture}[x=.85cm,y=.85cm,>=Stealth,font=\small]
 \path[use as bounding box] (-2.6,-2) rectangle (12.3,2.15);
 % Both panels use exactly the same coordinate scale and ellipse.
 \begin{scope}
  \node at (1,1.8) {外点：两个实切点};
  \draw[->,gray] (-2.3,0)--(4.5,0) node[right] {$x$};
  \draw[CStatement,thick] (0,0) ellipse (2 and 1);
  \draw[orange!80!black,thick] (1,.8660254)--(4,0)--(1,-.8660254);
  \draw[teal!80!black,very thick] (1,-1.25)--(1,1.25);
  \fill (4,0) circle(2pt) node[below] {$P(4,0)$};
  \fill (1,.8660254) circle(2pt); \fill (1,-.8660254) circle(2pt);
  \node[teal!80!black] at (.7,-1.6) {极线 $x=1$};
 \end{scope}
 \begin{scope}[shift={(7.5,0)}]
  \node at (1,1.8) {内点：没有实切点};
  \draw[->,gray] (-2.3,0)--(4.5,0) node[right] {$x$};
  \draw[CStatement,thick] (0,0) ellipse (2 and 1);
  \draw[teal!80!black,very thick] (4,-1.25)--(4,1.25);
  \fill (1,0) circle(2pt) node[below] {$P(1,0)$};
  \node[teal!80!black] at (3,-1.6) {极线 $x=4$};
 \end{scope}
\end{tikzpicture}

\end{center}

\textbf{错误想法：}写出极线，就必有两个实切点。
\textbf{看图纠正：}左图有两切点，右图极线无实交点；两图的椭圆与尺度相同。
先联立检查，有两个不同实交点才可称为切点弦。
\end{minipage}
\par


\medskip
\textbf{\textcolor{CTrap}{易错二｜忘记中心点没有普通极线}}
\begin{itemize}[leftmargin=*,itemsep=0.35em]
 \item 标准椭圆或双曲线的中心 $O(0,0)$ 对应 $0=1$，\textbf{它不是一条直线}。
 \item 对一般式，极线是直线须满足
 \[
 (2Ax_0+By_0+D,\ Bx_0+2Cy_0+E)\ne(0,0).
 \]
 \item 计算出只含常数的等式时，应停止，不要硬写成“直线方程”。
\end{itemize}

\medskip
\textbf{\textcolor{CTrap}{易错三｜$xy$ 项与一次项漏除以 $2$}}
\[
 \begin{array}{c|c|c}
 \text{原项}&\text{正确替换}&\text{常见错误}\\
 Bxy &\frac B2(x_0y+y_0x)&Bx_0y\\
 Dx &\frac D2(x+x_0)&Dx\\
 Ey &\frac E2(y+y_0)&Ey
 \end{array}
\]
\textbf{自检：}令 $x=x_0,y=y_0$，替换式必须还原为原曲线方程值。

\medskip
\textbf{\textcolor{CTrap}{易错四｜混淆椭圆、双曲线的符号}}
\[
 \text{椭圆：}\ \frac{x_0x}{a^2}+\frac{y_0y}{b^2}=1;
 \qquad
 \text{双曲线：}\ \frac{x_0x}{a^2}-\frac{y_0y}{b^2}=1.
\]
负号来自双曲线原方程本身，不能因记住椭圆公式就把它丢掉。

\medskip
\textbf{\textcolor{CTrap}{易错五｜只看到一个交点便称其为两切点连线}}
回看例 4 的图：极线 $x-y=1$ 只交曲线于 $T(1,0)$，橙色 $x=1$ 仍是切线。
\textbf{一个实切点既不等于没有切线，也不构成两切点连线。}

\medskip
\textbf{\textcolor{CTrap}{易错六｜误把任何二次方程都当作圆锥曲线}}
例如 $x^2-y^2=0$ 表示两条直线 $x=\pm y$。
使用前须确认方程确实表示椭圆、双曲线或抛物线。

\medskip
\textbf{\textcolor{CTrap}{考场检查顺序}}
\textbf{确认曲线种类 $\to$ 写出替换式 $\to$ 判断是否真是一条直线 $\to$
若题目要求切点弦，再联立检查是否有两个不同实切点。}
\end{trapbox}
